\lecture{Лекция 1}{10.09.2026}

\subsection{Полная система уравнений линейной теории упругости}

\begin{idea}[title={Обозначения}]
В тетради вектор помечен одной чертой снизу, тензор 2-го ранга — двумя.
В конспекте и то и другое набрано полужирным: $\vc u$ — вектор перемещения,
$\sig$ — тензор напряжений, $\Eten$ — единичный тензор. Двойное скалярное
произведение обозначено $\ddotp$\,.
\end{idea}

\subsubsection*{Геометрические соотношения}

Деформация определяется через симметричную часть градиента перемещения:
\begin{equation}\label{eq:geom}
  \epsT = \sym{\nab \vc u},
\end{equation}
где $\vc u$ — перемещение.

\subsubsection*{Уравнение равновесия}

\begin{equation}\label{eq:equilibrium}
  \divg\, \sig + \rho\left(\vc K - \dt{\vc v}\right) = 0 ,
\end{equation}
где $\vc K$ — массовая сила, $\dt{\vc v}$ — ускорение, $\rho$ — плотность.

\subsubsection*{Определяющие соотношения}

В простейшем случае (одномерное растяжение) закон Гука имеет вид $\sigma = E\eps$.
В общем случае напряжение — отклик материала на целый набор параметров:
\[
  \sig = \sig\left\{ T,\ \vc\Gamma,\ \tn F,\ \tn B \right\} ,
\]
где $T$ — температура, $\vc\Gamma$ — градиент температуры, $\tn F$ — градиент
деформации, $\tn B$ — тензор поворота. Аналогично устроены и остальные отклики:
\[
  \vc h = \dots \quad(\text{тепловой поток}), \qquad
  s = \dots \quad(\text{энтропия}).
\]

\begin{warn}
Список аргументов $\{T,\vc\Gamma,\tn F,\tn B\}$ и расшифровки $\tn F$, $\tn B$
взяты с доски: подписаны только $\vc\Gamma$ («град. температуры») и $\tn B$
(«тензор поворота»). Правые части для $\vc h$ и $s$ на лекции оставлены
многоточием.
\end{warn}

Дальше последовательно переходим от общего к частному. В самом простом случае
определяющее соотношение изотропного материала:
\begin{equation}\label{eq:hooke-K}
  \sig = K\theta\,\Eten + 2\mu\,\dev ,
  \qquad
  \sig = \lambda\theta\,\Eten + 2\mu\,\epsT .
\end{equation}

\subsubsection*{Учёт температуры}

Если появляется температура:
\begin{equation}\label{eq:hooke-temp}
  \sig = K\bigl(\theta - 3\alpha(T - T_0)\bigr)\Eten + 2\mu\,\dev ,
\end{equation}
где $T_0$ — температура, при которой деформаций нет, $\alpha$ — коэффициент
линейного температурного расширения. Первое слагаемое отвечает за
\textbf{объёмное расширение}, второе — за \textbf{изменение формы}.

При увеличении температуры тело расширяется — поэтому температура входит именно
в объёмную составляющую.

\begin{warn}
В тетради рядом стоит вопрос: «почему при $T$ стоит $-$?». Ответ на лекции не
записан.
\end{warn}

\begin{idea}
Знак минус нужен, чтобы при свободном нагреве напряжений не возникало. Свободное
тело при нагреве на $\Delta T = T - T_0$ получает объёмную деформацию
$\theta = 3\alpha\Delta T$, и скобка $\theta - 3\alpha(T-T_0)$ обращается в нуль:
шаровая часть напряжений равна нулю. Напряжение появляется только тогда, когда
расширению мешают — то есть реальная $\theta$ отличается от свободной
температурной.
\end{idea}

\subsubsection*{Шаровая часть и девиатор}

Разложение тензора деформации:
\[
  \epsT = \frac{1}{3} I_1(\epsT)\,\Eten + \dev ,
  \qquad
  \theta = I_1(\epsT) = \tr \epsT ,
\]
\[
  I_1(\epsT) = \eps_{11} + \eps_{22} + \eps_{33} .
\]
Величина
\[
  \theta = \frac{dv}{dV} - 1
\]
— \textbf{объёмная деформация} (дилатация): относительное изменение объёма,
$dV$ — объём в отсчётной конфигурации, $dv$ — в актуальной.

При нагреве, если ничто не мешает расширению, деформация равна
\[
  \theta = 3\alpha\,\Delta T .
\]

\subsubsection*{Граничные условия}

\inputfig{figures/fig-configurations}

Полная система в \textbf{эйлеровом подходе} (все поля заданы в актуальной
конфигурации, $\epsT = \epsT(\vc r)$):
\[
  \begin{cases}
    \epsT = \sym{\nab \vc u}, \\[0.3em]
    \divg\,\sig + \rho \vc K = 0, \\[0.3em]
    \sig = K\theta\,\Eten + 2\mu\,\dev .
  \end{cases}
\]

Варианты граничных условий.

\medskip
\textbf{1) Кинематические.} Задано перемещение на границе:
\[
  \vc u(\vc r)\big|_S = \vc U(\vc r) .
\]
Так как $\vc r = \vc R + \vc u$, разложение по Тейлору даёт
\[
  \vc u(\vc r)\big|_S = \vc u(\vc R + \vc u)\big|_S
  = \vc u(\vc R)\big|_S + \underbrace{\vc u \cdot (\nab \vc u)\big|_S}_{\to\,0}
  \approx \vc U(\vc R)\big|_S ,
\]
поскольку подчёркнутое слагаемое есть бесконечно малая большего порядка
при
$\norm{\nab \vc u} \ll 1$, $\norm{\vc u} \ll 1$. Точно так же
\[
  \vc U(\vc r) = \vc U(\vc R + \underbrace{\vc u}_{\text{беск. мало}}) .
\]
Отсюда
\[
  \boxed{\ \vc u(\vc R)\big|_S = \vc U(\vc R)\ }
\]
— условие можно сносить с актуальной конфигурации на отсчётную.

\medskip
\textbf{2) Статические.} Задана поверхностная сила:
\[
  \vc n \cdot \sig \big|_S = \vc F(\vc r)
  \quad\xrightarrow{\ \text{тем же рассуждением, что и в кин.}\ }\quad
  \vc n \cdot \sig \big|_S = \vc F(\vc R) .
\]

\medskip
\textbf{3) Смешанные.}
\[
  \text{на } S_1:\quad \vc u = \vc U ,
  \qquad
  \text{на } S_2:\quad \vc n \cdot \sig = \vc F .
\]

\begin{idea}[title={Счёт уравнений}]
Всего 15 уравнений: 6 геометрических соотношений~\eqref{eq:geom}, 3 уравнения
равновесия~\eqref{eq:equilibrium} и 6 определяющих
соотношений~\eqref{eq:hooke-K} — на 15 неизвестных ($3$ компоненты $\vc u$,
$6$ компонент $\epsT$, $6$ компонент $\sig$). Плюс 3 граничных условия в каждой
точке границы.
\end{idea}
